% =============================================================
% บทที่ 3 - วิธีการดำเนินโครงงาน
% =============================================================

% -----------------------------------------------------------
\section{ขั้นตอนการดำเนินงาน}
% -----------------------------------------------------------

โครงงานนี้ดำเนินงานตามหลักการวิศวกรรมซอฟต์แวร์ โดยแบ่งขั้นตอนการดำเนินงานเป็น 5 ขั้นตอนหลัก ดังนี้

\subsection{ศึกษาปัญหาและรวบรวมความต้องการ}

ศึกษาและวิเคราะห์ปัญหาในการติดตามข่าวสารจากหลายสำนักข่าวไทย รวบรวมความต้องการเชิงฟังก์ชันและไม่ใช่เชิงฟังก์ชันของระบบ ศึกษาโครงสร้างเว็บไซต์และ RSS Feed ของสำนักข่าวเป้าหมาย กำหนดชุดหมวดหมู่ข่าว และเลือกใช้โมเดลภาษาขนาดใหญ่ผ่านบริการ KKU Gen.AI

\subsection{วิเคราะห์และออกแบบระบบ}

วิเคราะห์ความต้องการและออกแบบสถาปัตยกรรมของระบบแบบแบ่งชั้น (Layered Architecture) ตามหลักการ SOLID และ GRASP ประกอบด้วยชั้น Routers, Services, Scrapers, Repository และ Core โดยกำหนดแผนภาพ Use Case Diagram สำหรับผู้ใช้งานและระบบ ออกแบบโครงสร้างข้อมูลข่าวด้วย Pydantic Model และกำหนดโปรโตคอล Repository ด้วย Python Protocol เพื่อความยืดหยุ่นในการเปลี่ยนวิธีจัดเก็บข้อมูล พร้อมทั้งออกแบบโครงสร้างอัลกอริทึมการกำหนดคุกกี้ (Assign Cookies Algorithm) เพื่อใช้เป็นกลไกหลักในการคำนวณความสนใจและปรับเปลี่ยนเนื้อหาให้สอดคล้องกับพฤติกรรมผู้ใช้แต่ละคน

\subsection{พัฒนาระบบ}

พัฒนาระบบตามที่ออกแบบไว้ โดยใช้ภาษา Python ร่วมกับเฟรมเวิร์ค FastAPI สำหรับฝั่ง Backend พัฒนา Scraper สำหรับแต่ละสำนักข่าวด้วยการลงทะเบียนแบบ Decorator พัฒนาบริการสรุปข่าวโดยส่งเนื้อหา Markdown ไปยังโมเดลภาษาขนาดใหญ่ พัฒนาบริการจำแนกหมวดหมู่ด้วย PyThaiNLP ร่วมกับโมเดล SVM และพัฒนา Frontend ด้วย JavaScript แบบ Vanilla ร่วมกับ Tailwind CSS และ Socket.IO สำหรับการแสดงผลแบบเรียลไทม์

นอกจากนี้ ยังพัฒนากลไกแคชเพื่อจัดเก็บข้อมูลข่าวและเนื้อหาที่ดึงมาแล้ว โดยใช้ URL เป็นกุญแจแคช เพื่อหลีกเลี่ยงการดึงข้อมูลซ้ำในแต่ละรอบการทำงาน และพัฒนาระบบจัดการคุกกี้สามประเภท ได้แก่ (1) คุกกี้เพื่อการดึงข้อมูล ร่วมกับ Playwright Service เพื่อคงสถานะเซสชันของเบราว์เซอร์สำหรับการดึงข้อมูลจากเว็บไซต์ที่แสดงผลเนื้อหาผ่าน JavaScript หรือมีกลไกตรวจสอบเบราว์เซอร์ (2) คุกกี้เพื่อการปรับแต่งเฉพาะบุคคล (Personalize Cookies) สำหรับจัดเก็บค่ากำหนดของผู้ใช้ เช่น หมวดหมู่ที่สนใจและการตั้งค่าการแสดงผล และ (3) คุกกี้เพื่ออัลกอริทึม (Algorithm Cookies) สำหรับบันทึกพฤติกรรมการอ่านเพื่อป้อนข้อมูลให้อัลกอริทึมแนะนำข่าว

\subsection{ทดสอบและประเมินผล}

ทดสอบแต่ละส่วนของระบบด้วยการตรวจสอบโค้ดด้วย Ruff และตรวจสอบชนิดข้อมูลด้วย mypy ทดสอบการทำงานของ Repository และ Service ด้วยหน่วยทดสอบ (Unit Test) ทดสอบการดึงข่าวของแต่ละสำนักข่าวว่าสามารถดึงข้อมูลได้ครบถ้วน และประเมินความถูกต้องของบทสรุปจากโมเดลภาษาขนาดใหญ่และความถูกต้องของการจำแนกหมวดหมู่ข่าว รวมถึงทดสอบการแจ้งเตือนข่าวแบบเรียลไทม์ผ่าน WebSocket

\subsection{ปรับปรุงและสรุปผล}

ปรับปรุงระบบตามผลการทดสอบ จัดทำเอกสารรายงานโครงงานและคู่มือการใช้งาน และสรุปผลการดำเนินงานเทียบกับวัตถุประสงค์ที่กำหนดไว้

% -----------------------------------------------------------
\section{เครื่องมือที่ใช้ในการดำเนินงาน}
% -----------------------------------------------------------

\subsection{ฮาร์ดแวร์}

เครื่องมือด้านฮาร์ดแวร์ที่ใช้ในการพัฒนาโครงงานแสดงดังตารางที่ \ref{tab:hardware}

\begin{table}[H]
\centering
\caption{รายการฮาร์ดแวร์ที่ใช้ในการพัฒนา}
\label{tab:hardware}
\begin{tabular}{|c|l|p{6.5cm}|}
\hline
\textbf{ลำดับ} & \textbf{รายการ} & \textbf{รายละเอียด} \\
\hline
1 & คอมพิวเตอร์ส่วนบุคคล & CPU Intel Core Ultra 5, RAM 16 GB, SSD 512 GB \\ \hline
2 & จอภาพ & ขนาด 16 นิ้ว \\ \hline
3 & การเชื่อมต่ออินเทอร์เน็ต & ความเร็วไม่น้อยกว่า 30 Mbps ใช้ดึงข่าวและเรียกใช้บริการ KKU Gen.AI \\
\hline
\end{tabular}
\end{table}

\subsection{ซอฟต์แวร์}

เครื่องมือด้านซอฟต์แวร์ที่ใช้ในการพัฒนาโครงงานแสดงดังตารางที่ \ref{tab:software}

\begin{table}[H]
\centering
\caption{รายการซอฟต์แวร์ที่ใช้ในการพัฒนา}
\label{tab:software}
\begin{tabular}{|c|l|c|p{5.5cm}|}
\hline
\textbf{ลำดับ} & \textbf{ซอฟต์แวร์} & \textbf{เวอร์ชัน} & \textbf{วัตถุประสงค์} \\
\hline
1 & Windows/Linux & \ldots & สภาพแวดล้อมการพัฒนา \\ \hline
2 & Python & 3.10 ขึ้นไป & พัฒนาระบบ Backend \\ \hline
3 & FastAPI + Uvicorn & \ldots & พัฒนา REST API และเว็บเซิร์ฟเวอร์ \\ \hline
4 & httpx + BeautifulSoup4 & \ldots & ดึงและวิเคราะห์หน้าเว็บ \\ \hline
5 & Trafilatura & \ldots & สกัดเนื้อหาบทความเป็น Markdown \\ \hline
6 & Playwright & \ldots & ดึงข้อมูลเว็บไซต์ที่ใช้ JavaScript \\ \hline
7 & PyThaiNLP & \ldots & ตัดคำภาษาไทยสำหรับจำแนกหมวดหมู่ \\ \hline
8 & scikit-learn, joblib & \ldots & โมเดล SVM/TF-IDF จำแนกหมวด \\ \hline
9 & OpenAI SDK & \ldots & เรียกใช้โมเดล LLM สรุปข่าว \\ \hline
10 & python-socketio & \ldots & แจ้งเตือนข่าวแบบเรียลไทม์ \\ \hline
11 & Node.js + Tailwind CSS & 18+ / \ldots & พัฒนา Frontend \\ \hline
12 & Git + GitHub & \ldots & จัดการเวอร์ชันของซอฟต์แวร์ \\ \hline
13 & Docker & \ldots & จัดการบริการ Playwright Service \\ \hline
14 & กลไกแคช (Cache) & \ldots & จัดเก็บข้อมูลข่าวที่ดึงแล้ว ลดการดึงข้อมูลซ้ำ และเพิ่มความเร็วในการแสดงผล \\ \hline
15 & ระบบคุกกี้และอัลกอริทึมแนะนำ & \ldots & จัดเก็บค่ากำหนดผู้ใช้ บันทึกพฤติกรรมการอ่าน และแนะนำข่าวเฉพาะบุคคล \\
\hline
\end{tabular}
\end{table}

% -----------------------------------------------------------
\section{แผนการดำเนินงาน}
% -----------------------------------------------------------

แผนการดำเนินงานตลอดระยะเวลาการทำโครงงาน 1 ภาคการศึกษา (6 เดือน) แสดงดังตารางที่ \ref{tab:gantt}

\begin{table}[H]
\centering
\caption{แผนการดำเนินงาน (Gantt Chart)}
\label{tab:gantt}
\begin{tabular}{|c|p{3.0cm}|c|c|c|c|c|c|}
\hline
\textbf{ลำดับ} & \textbf{กิจกรรม} & \textbf{เดือน 1} & \textbf{เดือน 2} & \textbf{เดือน 3} & \textbf{เดือน 4} & \textbf{เดือน 5} & \textbf{เดือน 6} \\ \hline
1 & ศึกษาปัญหา/ความต้องการ & $\bullet$ & & & & & \\ \hline
2 & วิเคราะห์และออกแบบ & & $\bullet$ & $\bullet$ & & & \\ \hline
3 & ออกแบบอัลกอริทึมกำหนดคุกกี้ & & $\bullet$ & & & & \\ \hline
4 & พัฒนา Backend/Scraper & & & $\bullet$ & $\bullet$ & & \\ \hline
5 & พัฒนา Frontend+Realtime & & & & $\bullet$ & $\bullet$ & \\ \hline
6 & ทดสอบและประเมินผล & & & & & $\bullet$ & $\bullet$ \\ \hline
7 & ปรับปรุงและจัดทำรายงาน & & & & & & $\bullet$ \\ \hline
\end{tabular}
\end{table}
